% Part: turing-machines
% Chapter: undecidability
% Section: unsolvability-decision-problem

\documentclass[../../../include/open-logic-section]{subfiles}

\begin{document}

\olfileid{tur}{und}{dec}
\olsection{નિર્ણય સમસ્યા}

આપેલું !!{sentence} પ્રામાણ્ય ધરાવે છે કે નહીં તે
નક્કી કરવાની અસરકારક રીત હોય ત્યારે અને માત્ર ત્યારે
જ આપણે પ્રથમ-ક્રમનું તર્કશાસ્ત્ર \emph{નિર્ણેય} છે એમ
કહીએ છીએ. હકીકતમાં, આવી કોઈ રીત નથી: પ્રથમ-ક્રમનાં
વાક્યોનું પ્રામાણ્ય નક્કી કરવાની સમસ્યા ઉકેલી શકાતી નથી.

આ મહત્ત્વનું નકારાત્મક પરિણામ સ્થાપિત કરવા આપણે
નિર્ણય સમસ્યા ટ્યુરિંગ મશીન વડે ઉકેલી શકાતી નથી તે
સાબિત કરીએ છીએ. એટલે, એવું કોઈ ટ્યુરિંગ મશીન નથી
જેની ટેપ પર પ્રથમ-ક્રમનું !!{sentence} હોય ત્યારે
તે શરૂ કરીને છેવટે અટકે અને તે !!{sentence} પ્રામાણ્ય ધરાવે
છે કે નહીં તે મુજબ $1$ અથવા~$0$ આઉટપુટ આપે.
ચર્ચ--ટ્યુરિંગ માન્યતા મુજબ દરેક સંગણનીય વિધેય
ટ્યુરિંગ-સંગણનીય છે. તેથી જો આ ``પ્રામાણ્ય વિધેય''
કોઈ પણ રીતે અસરકારક રીતે સંગણનીય હોત, તો તે
ટ્યુરિંગ-સંગણનીય હોત. તે ટ્યુરિંગ-સંગણનીય નથી,
તો અસરકારક રીતે પણ સંગણનીય હોઈ શકે નહીં.

નિર્ણય સમસ્યા ઉકેલી શકાતી નથી તે સાબિત કરવાની
આપણી રીત અટકવાની સમસ્યાનું તેમાં ન્યૂનીકરણ કરવાની
છે. તેનો અર્થ આ છે: ટ્યુરિંગ મશીનનો સૂચક~$e$
અને ઇનપુટ~$w$ માટે, મશીન અટકે તો $1$ અને
અન્યથા~$0$ મૂલ્ય આપતું વિધેય~$h(e,w)$ ટ્યુરિંગ-સંગણનીય
નથી એ આપણે સાબિત કર્યું છે. આપણે બતાવીશું કે
પ્રથમ-ક્રમનાં વાક્યોનું પ્રામાણ્ય નક્કી કરતું ટ્યુરિંગ
મશીન હોત, તો~$h$ ગણતું ટ્યુરિંગ મશીન પણ હોત.
પરંતુ~$h$ને ટ્યુરિંગ મશીન વડે ગણી શકાતું નથી;
તેથી પ્રામાણ્ય નક્કી કરતું મશીન પણ હોઈ શકે નહીં.

આ રીતનું પહેલું પગલું એ બતાવવાનું છે કે દરેક ઇનપુટ~$w$
અને ટ્યુરિંગ મશીન~$M$ માટે આપણે અસરકારક રીતે
!!a{sentence} $!T(M, w)$ આપી શકીએ, જે~$M$ની
સૂચનાઓનો ગણ અને ઇનપુટ~$w$ નિરૂપે, અને
!!a{sentence}~$!E(M, w)$ આપી શકીએ, જે ``$M$ છેવટે
અટકે છે'' એવું વ્યક્ત કરે, એવી રીતે કે:
\begin{quote}
  $\Entails !T(M, w) \lif !E(M,w)$ ત્યારે અને માત્ર
  ત્યારે જ, જ્યારે $M$ ઇનપુટ~$w$ પર અટકે છે.
\end{quote}
આ વાક્યો $!T(M, w)$ અને~$!E(M, w)$નું વર્ણન અને
$!T(M, w) \lif !E(M, w)$ ત્યારે અને માત્ર ત્યારે જ
પ્રામાણ્ય ધરાવે છે જ્યારે $M$ ઇનપુટ~$w$ પર અટકે
છે તેની તપાસ, આપણી સાબિતીનો મુખ્ય ભાગ હશે.

\end{document}
